\documentclass[12pt, a4paper]{article}

% ── Packages ──────────────────────────────────────────────────────
\usepackage[utf8]{inputenc}
\usepackage[T1]{fontenc}
\usepackage{lmodern}
\usepackage[margin=1in]{geometry}
\usepackage{graphicx}
\usepackage{booktabs}
\usepackage{hyperref}
\usepackage{xcolor}
\usepackage{amsmath}
\usepackage{caption}
\usepackage{subcaption}
\usepackage{float}
\usepackage{enumitem}
\usepackage{parskip}
\usepackage{fancyhdr}
\usepackage{titlesec}
\usepackage{multirow}

\hypersetup{
  colorlinks=true,
  linkcolor=blue!60!black,
  citecolor=blue!60!black,
  urlcolor=blue!60!black,
}

% ── Headers ───────────────────────────────────────────────────────
\pagestyle{fancy}
\fancyhf{}
\fancyhead[L]{\small Pearls AQI Predictor}
\fancyhead[R]{\small \thepage}
\renewcommand{\headrulewidth}{0.4pt}

% ── Title ─────────────────────────────────────────────────────────
\title{
  \textbf{Pearls AQI Predictor} \\[0.3em]
  \large Serverless Multi-Horizon Air Quality Forecasting for Karachi \\[0.5em]
  \normalsize End-to-End ML Pipeline Report
}
\author{Abdul}
\date{September 2026}

% ══════════════════════════════════════════════════════════════════
\begin{document}
\maketitle
\thispagestyle{empty}

\begin{abstract}
This report documents the design, implementation, and evaluation of Pearls AQI Predictor --- a serverless machine learning system that forecasts the US Air Quality Index (AQI) for Karachi up to 72 hours ahead.  The system ingests hourly weather and pollutant data from the Open-Meteo API, engineers leakage-free features, trains direct per-horizon models (XGBoost, Random Forest, LSTM), and serves real-time predictions through a FastAPI backend and React dashboard.  Automation is handled by GitHub Actions (hourly feature updates, daily retraining) and deployment is managed through FastAPI Cloud and Vercel.  The best model (XGBoost) achieves an RMSE of 0.92 at the +1h horizon ($R^2 = 0.995$) and 9.41 at +24h ($R^2 = 0.511$).  SHAP analysis confirms the model relies on physically meaningful features and is free of target leakage.
\end{abstract}

\tableofcontents
\newpage

% ══════════════════════════════════════════════════════════════════
\section{Introduction}

Air pollution is a serious public health concern in Karachi, Pakistan's largest city.  The US Air Quality Index (AQI) condenses concentrations of multiple pollutants into a single number that communicates health risk to the public.  This project builds a complete, serverless pipeline that:

\begin{enumerate}
  \item Collects hourly weather and pollutant data from a free public API.
  \item Engineers time-aware features and stores them in a cloud feature store.
  \item Trains separate ML models for six forecast horizons (1h, 6h, 12h, 24h, 48h, 72h).
  \item Serves forecasts through a REST API and an interactive web dashboard.
  \item Automates the entire cycle with GitHub Actions.
\end{enumerate}

The goal is to predict the AQI for the \emph{next three days} with enough accuracy to be useful for health advisories, while keeping the infrastructure entirely serverless and free-tier friendly.


% ══════════════════════════════════════════════════════════════════
\section{Data Sources and Collection}

\subsection{Open-Meteo API}

All data comes from the \textbf{Open-Meteo API} (\texttt{open-meteo.com}), which provides free, no-key-required access to both weather and air quality data.  Open-Meteo was chosen over AQICN because it offers both weather \emph{and} air quality from a single provider, provides a forecast API (necessary for multi-day predictions), and has no rate-limit issues for hourly automated pipelines.

Two types of hourly variables are fetched:

\begin{itemize}
  \item \textbf{Weather variables} (6): temperature, relative humidity, precipitation, sea-level pressure, wind speed, wind direction.
  \item \textbf{Air quality variables} (7): PM\textsubscript{2.5}, PM\textsubscript{10}, CO, NO\textsubscript{2}, SO\textsubscript{2}, O\textsubscript{3}, and the US AQI (the prediction target).
\end{itemize}

\subsection{Location}

The system is configured for \textbf{Karachi} (24.86\textdegree N, 67.01\textdegree E).  The city was chosen because it has consistently poor air quality (mean AQI $\approx 88$), making forecasting both challenging and practically useful.  The location is configurable through environment variables.

\subsection{Historical Backfill}

A one-time backfill pipeline (\texttt{backfill.py}) was run to populate the feature store with approximately three years of historical data (September 2023 -- September 2026, totaling 26,496 hourly rows).  The Open-Meteo historical API limits data per request, so the backfill processes the date range in 90-day chunks.  Each chunk after the first includes a 4-day overlap buffer so that lag and rolling features (up to 72h) compute correctly at chunk boundaries.


% ══════════════════════════════════════════════════════════════════
\section{Exploratory Data Analysis}

An EDA was conducted in two phases: a notebook-based initial analysis (\texttt{01\_eda\_open\_meteo\_aqi.ipynb}) covering distributions, correlations, and basic patterns; and a script-based advanced analysis (\texttt{eda\_advanced.py}) covering stationarity testing, ACF/PACF, seasonal decomposition, and outlier characterization.

\subsection{Target Variable Distribution}

The US AQI for Karachi has the following summary statistics:

\begin{table}[H]
\centering
\caption{US AQI descriptive statistics (2-year raw dataset, 17,544 rows).}
\begin{tabular}{lr}
\toprule
\textbf{Statistic} & \textbf{Value} \\
\midrule
Mean   & 88.0 \\
Std    & 22.0 \\
Min    & 41 \\
Median & 82 \\
Max    & 173 \\
\bottomrule
\end{tabular}
\end{table}

The distribution is right-skewed, with the majority of readings falling in the ``Moderate'' EPA band (51--100 AQI) and a long tail reaching into ``Unhealthy'' (151--200).  No readings cross the ``Very Unhealthy'' (201--300) or ``Hazardous'' ($>$300) thresholds in this dataset.

\subsection{Temporal Patterns}

\begin{itemize}
  \item \textbf{Seasonal cycle}: AQI is markedly higher and more volatile during winter (October--February) due to temperature inversions that trap pollutants.  Summer monsoon rains wash out particulates, producing lower and calmer AQI values.
  \item \textbf{Diurnal cycle}: AQI peaks around 7--8 PM ($\sim$96 AQI) and is lowest in the early morning.  This aligns with evening rush-hour emissions coinciding with the collapse of the atmospheric boundary layer at sunset.
  \item \textbf{Day-of-week}: No meaningful weekly pattern was found (all days within 88--91 AQI), so day-of-week was kept only as a weak secondary feature.
\end{itemize}

\subsection{Stationarity and Autocorrelation}

The Augmented Dickey-Fuller (ADF) test confirmed that the raw AQI series is \textbf{stationary} ($p \ll 0.05$), meaning no differencing is needed.

ACF analysis revealed significant autocorrelation at all lags up to 72 hours (ACF = 0.440 at lag 72h), confirming that past AQI values are informative even for 3-day-ahead forecasts.  PACF showed direct influence at 24h, 48h, and 72h lags, indicating a strong diurnal cycle.  These findings directly informed the choice of lag horizons in the feature engineering step.

\subsection{Correlation Analysis}

The strongest linear correlates with AQI are:
\begin{itemize}
  \item 6h rolling mean AQI (correlation = 0.96) --- the dominant predictor for short horizons.
  \item PM\textsubscript{2.5} (0.72) --- the primary pollutant driver of AQI.
  \item Wind speed ($-$0.41) --- higher wind disperses pollutants.
  \item Temperature ($-$0.35) and relative humidity ($-$0.33) --- warmer, humid monsoon conditions correspond to cleaner air.
\end{itemize}

Hour and month showed near-zero \emph{linear} correlations despite clear visual patterns, because their effects are cyclical.  This motivated the use of sin/cos encoding rather than raw integer features.


% ══════════════════════════════════════════════════════════════════
\section{Feature Engineering}
\label{sec:features}

Feature engineering is handled by \texttt{build\_features.py}, which transforms the 13 raw columns into approximately 60 engineered features.  The features are grouped into six categories:

\subsection{Time-Based Features}

Hours, months, and days of the week are encoded as \textbf{sin/cos pairs} (e.g., $\sin(2\pi \cdot \text{hour}/24)$) so that the model understands that hour 23 is close to hour 0.  A binary weekend flag is also included.

\subsection{AQI Lag Features}

Lagged values of the target variable at 1h, 3h, 6h, 12h, 24h, 48h, and 72h.  The 48h and 72h lags were added based on the ACF/PACF analysis showing significant autocorrelation at those horizons.

\subsection{AQI Rolling and Statistical Features}

Rolling means (3h, 6h, 12h, 24h, 48h), rolling standard deviation (6h, 24h), rolling min/max (24h), and an exponentially weighted mean with 12h half-life.  All rolling computations are \textbf{shifted by 1 hour} to prevent leaking the current value into the window.

\subsection{Pollutant Lag and Rolling Features}

Lagged and rolling averages for PM\textsubscript{2.5}, PM\textsubscript{10}, CO, NO\textsubscript{2}, SO\textsubscript{2}, and O\textsubscript{3}.  Raw pollutant values at the current timestep are \textbf{dropped} to prevent target leakage, since the US AQI is a deterministic function of these pollutant concentrations.  Only their \emph{past} values (lags) and \emph{historical averages} (rolling means) are retained.

\subsection{Trend and Change Features}

First differences of AQI at 1h, 3h, 6h, and 24h intervals capture the rate of change, and deviation from rolling means (trend features) captures whether air quality is currently improving or worsening.

\subsection{Weather Interaction Features}

\begin{itemize}
  \item \textbf{Wind vector decomposition}: Wind speed and direction are converted to u/v vector components for the model to reason about wind direction continuously.
  \item \textbf{Temperature-humidity index}: The product $T \times H / 100$ serves as a proxy for conditions that favor secondary particulate formation.
  \item \textbf{Precipitation indicators}: Cumulative rain in the last 6h and a binary flag for any rain in the last 24h.
\end{itemize}


% ══════════════════════════════════════════════════════════════════
\section{Feature Store}

Processed features are stored in \textbf{Hopsworks Feature Store} (cloud-hosted, free tier).  Hopsworks was chosen because it provides managed storage with time-travel support (HUDI format), a Python SDK for easy integration, and a model registry for trained models.

The feature group (\texttt{aqi\_features}, version 2) uses \texttt{time} as both the primary key and the event-time column.  Version 2 represents the leakage-free feature set; version 1 (which included raw pollutant columns) is retained for reference but not used.

An insert-with-retry mechanism handles transient network errors during backfill and hourly updates.


% ══════════════════════════════════════════════════════════════════
\section{Modeling Approach}
\label{sec:modeling}

\subsection{Direct Multi-Horizon Strategy}

Rather than training a single +1h model and feeding its predictions back recursively, the system trains \textbf{a separate model for each forecast horizon}: +1h, +6h, +12h, +24h, +48h, and +72h.  This ``direct'' approach was chosen because:

\begin{enumerate}
  \item Recursive forecasting accumulates prediction errors at each step, which becomes severe at 72 steps.
  \item Each horizon can use a feature set tailored to what information is actually useful at that distance.
  \item Training is embarrassingly parallel --- each horizon model is independent.
\end{enumerate}

For intermediate hours (e.g., +2h through +5h), predictions are linearly interpolated between the nearest direct-model outputs.

\subsection{Forecast Weather Integration}

A key insight is that Open-Meteo provides \textbf{3-day weather forecasts}, and this future weather data is legitimately available at inference time.  The training pipeline simulates this by shifting weather columns by $-h$ hours to create ``forecast weather'' features for each horizon.  Additionally, \textbf{weather deltas} (forecast minus current) are computed, capturing the expected meteorological change between now and the target time.

For example, the +24h model receives both the current temperature \emph{and} the forecasted temperature 24 hours from now, plus the difference between them.

\subsection{Horizon-Specific Feature Selection}

For short horizons ($\leq$6h), all features are used.  For long horizons ($\geq$24h), short-term features like 1h and 3h lags are pruned because they will have decayed into noise by the target time.  The model is instead anchored on 24h diurnal lags, rolling means, and forecast weather deltas.  This prevents the model from overfitting to features that carry no signal at the prediction horizon.

\subsection{Models Evaluated}

Three model families were trained for each horizon:

\begin{enumerate}
  \item \textbf{XGBoost} (gradient-boosted trees): 500 estimators, max depth 6, learning rate 0.05, with early stopping on a validation set to prevent overfitting.
  \item \textbf{Random Forest}: 300 trees, max depth 15, used as a stable baseline.
  \item \textbf{LSTM} (PyTorch): 2-layer LSTM with 128 hidden units, 48-hour sliding window, Huber loss ($\delta=10$), and learning rate scheduling.  The LSTM uses StandardScaler-transformed inputs and scaled targets.
\end{enumerate}

The data split is chronological --- 70\% train, 15\% validation, 15\% test --- with no shuffling, to respect the temporal ordering.


% ══════════════════════════════════════════════════════════════════
\section{Results}

\subsection{Model Performance Comparison}

Table~\ref{tab:results} shows the test-set performance of all three models across the six horizons.  XGBoost is the best model at every horizon.

\begin{table}[H]
\centering
\caption{Test-set metrics for all models and horizons.  Bold indicates the best model per horizon.}
\label{tab:results}
\small
\begin{tabular}{llrrr}
\toprule
\textbf{Horizon} & \textbf{Model} & \textbf{RMSE} & \textbf{MAE} & \textbf{$R^2$} \\
\midrule
\multirow{3}{*}{+1h}
  & \textbf{XGBoost}     & \textbf{0.92} & \textbf{0.46} & \textbf{0.995} \\
  & Random Forest        & 1.12          & 0.48          & 0.993 \\
  & LSTM                 & 1.73          & 0.82          & 0.984 \\
\midrule
\multirow{3}{*}{+6h}
  & \textbf{XGBoost}     & \textbf{3.94} & \textbf{2.03} & \textbf{0.914} \\
  & Random Forest        & 4.74          & 2.28          & 0.876 \\
  & LSTM                 & 4.89          & 2.50          & 0.869 \\
\midrule
\multirow{3}{*}{+12h}
  & \textbf{XGBoost}     & \textbf{6.34} & \textbf{3.57} & \textbf{0.778} \\
  & Random Forest        & 6.76          & 3.78          & 0.748 \\
  & LSTM                 & 7.05          & 4.09          & 0.728 \\
\midrule
\multirow{3}{*}{+24h}
  & \textbf{XGBoost}     & \textbf{9.41} & \textbf{6.25} & \textbf{0.511} \\
  & Random Forest        & 9.61          & 6.38          & 0.490 \\
  & LSTM                 & 9.82          & 6.69          & 0.472 \\
\midrule
\multirow{3}{*}{+48h}
  & \textbf{XGBoost}     & \textbf{12.07} & \textbf{8.29} & \textbf{0.200} \\
  & Random Forest        & 12.28          & 8.49          & 0.172 \\
  & LSTM                 & 14.26          & 9.83          & $-$0.136 \\
\midrule
\multirow{3}{*}{+72h}
  & \textbf{XGBoost}     & \textbf{12.60} & \textbf{8.90} & \textbf{0.131} \\
  & Random Forest        & 13.25          & 9.25          & 0.040 \\
  & LSTM                 & 13.00          & 9.09          & 0.042 \\
\bottomrule
\end{tabular}
\end{table}

\subsection{Discussion of Results}

\begin{itemize}
  \item \textbf{+1h}: Near-perfect accuracy (RMSE $<$ 1 AQI point, $R^2 = 0.995$).  The model essentially learns that AQI rarely changes by more than 1--2 points in an hour, and the 1h lag feature dominates.
  \item \textbf{+6h to +12h}: Strong performance.  The model relies on rolling means and pollutant lag features.  Forecast weather deltas begin to add value.
  \item \textbf{+24h}: Moderate performance ($R^2 = 0.51$).  At this horizon, the model shifts to diurnal lags and forecast weather, consistent with the AQI being driven by the daily meteorological cycle.
  \item \textbf{+48h and +72h}: Performance degrades significantly ($R^2 = 0.13$--$0.20$).  This is expected --- 2--3 day AQI forecasting is inherently difficult because pollution levels depend on future emission events and mesoscale meteorology that are only weakly predictable.  The LSTM performs especially poorly at +48h ($R^2 < 0$, worse than predicting the mean).
\end{itemize}

The fact that RMSE only increases modestly from +48h (12.07) to +72h (12.60) suggests the model is not adding much signal beyond the 48h horizon and is mostly returning to the climatological mean.

\subsection{Comparison with Baselines}

Naive baselines were calculated to calibrate model performance:

\begin{itemize}
  \item \textbf{Mean baseline}: Always predict the training-set mean AQI. By definition, $R^2 = 0$.
  \item \textbf{Persistence}: Predict $\text{AQI}(t+h) = \text{AQI}(t)$ (``it stays the same'').
  \item \textbf{Same-hour-yesterday}: Predict $\text{AQI}(t+h) = \text{AQI}(t+h-24)$.
\end{itemize}

The XGBoost models beat all three baselines at every horizon, confirming that the learned models provide genuine predictive value beyond simple heuristics.


% ══════════════════════════════════════════════════════════════════
\section{Explainability with SHAP}

SHAP (SHapley Additive exPlanations) analysis was run on the XGBoost models for the +1h and +24h horizons to understand what drives short-term vs.\ long-term predictions.

\subsection{+1h Model}

The top features by mean absolute SHAP value are:

\begin{table}[H]
\centering
\caption{Top 10 features for the +1h XGBoost model by SHAP importance.}
\begin{tabular}{lr}
\toprule
\textbf{Feature} & \textbf{Mean $|$SHAP$|$} \\
\midrule
us\_aqi\_lag\_1h       & 9.175 \\
pm2\_5\_rolling\_24h   & 3.325 \\
ozone\_rolling\_6h     & 1.067 \\
aqi\_change\_1h        & 0.629 \\
aqi\_rolling\_3h       & 0.447 \\
pm2\_5\_lag\_1h         & 0.254 \\
pm2\_5\_rolling\_12h   & 0.253 \\
aqi\_rolling\_6h       & 0.191 \\
pm10\_rolling\_6h      & 0.168 \\
aqi\_ewm\_12h          & 0.131 \\
\bottomrule
\end{tabular}
\end{table}

The 1h AQI lag dominates by a wide margin, which is physically intuitive: AQI one hour from now is mostly determined by AQI right now.  The second most important feature is the 24h rolling mean of PM\textsubscript{2.5}, which provides a stable baseline of pollution levels.

\textbf{Leakage check}: No raw pollutant columns (pm2\_5, pm10, etc.) appear in the model --- only their lagged and rolling derivatives.  This confirms the leakage fix is working correctly.

\subsection{+24h Model}

At the +24h horizon, the SHAP importance shifts substantially:

\begin{table}[H]
\centering
\caption{Top 7 features for the +24h XGBoost model by SHAP importance.}
\begin{tabular}{lr}
\toprule
\textbf{Feature} & \textbf{Mean $|$SHAP$|$} \\
\midrule
pm2\_5\_rolling\_6h           & 4.458 \\
month\_sin                    & 1.984 \\
pm2\_5\_rolling\_24h          & 1.870 \\
pm2\_5\_rolling\_12h          & 1.814 \\
forecast\_temp\_humidity\_24h & 1.437 \\
aqi\_change\_6h               & 1.185 \\
ozone\_rolling\_6h            & 0.948 \\
\bottomrule
\end{tabular}
\end{table}

Key differences from the +1h model:
\begin{itemize}
  \item The 1h AQI lag (dominant at +1h) is \emph{pruned} from the +24h model because it carries no useful signal 24 hours ahead.
  \item \textbf{Month} (via sin encoding) becomes the second most important feature, reflecting the strong seasonal AQI pattern (higher in winter, lower in monsoon).
  \item \textbf{Forecast temperature-humidity index} (\texttt{forecast\_temp\_humidity\_24h}) appears, showing the model uses future weather conditions to adjust its prediction.
  \item PM\textsubscript{2.5} rolling averages across multiple windows (6h, 12h, 24h) collectively dominate, providing a multi-scale view of the prevailing pollution baseline.
\end{itemize}


% ══════════════════════════════════════════════════════════════════
\section{System Architecture and Deployment}

The system consists of four layers:

\subsection{Data Layer}

\texttt{sources.py} handles all API calls to Open-Meteo.  It merges weather and air quality data on the time column and returns a clean DataFrame.

\subsection{Feature Store Layer}

\texttt{hopsworks\_utils.py} manages the connection to Hopsworks, creates or retrieves the feature group, and inserts data with retry logic (up to 3 attempts with exponential backoff).

\subsection{Serving Layer (Backend)}

A \textbf{FastAPI} application (\texttt{api.py}) provides four main endpoints:

\begin{itemize}
  \item \texttt{/api/current} --- Returns the latest AQI reading with EPA category, dominant pollutant, and weather conditions.
  \item \texttt{/api/forecast} --- Generates a 72-hour forecast using the direct XGBoost models, with per-hour AQI predictions and weather context.
  \item \texttt{/api/history} --- Returns the last 48 hours of observed AQI data.
  \item \texttt{/api/analytics} --- Returns model evaluation metrics and SHAP feature importance for the dashboard's explainability panel.
\end{itemize}

The API implements a 5-minute TTL cache for weather data to avoid hitting the Open-Meteo API on every request.  The backend is deployed on \textbf{FastAPI Cloud}, which auto-redeploys when model artifacts are pushed to the repository.

\subsection{Frontend (Dashboard)}

The dashboard is a \textbf{React + TypeScript} application built with Vite, styled with Tailwind CSS v4, and using Recharts for data visualization.  It includes:

\begin{itemize}
  \item A live AQI hero card with EPA category color coding and health advisory.
  \item Pollutant concentration cards (PM\textsubscript{2.5}, PM\textsubscript{10}, CO, NO\textsubscript{2}, SO\textsubscript{2}, O\textsubscript{3}).
  \item A 72-hour interactive forecast chart with daily summary cards.
  \item A 48-hour historical observation chart.
  \item A model analytics and SHAP explainability panel showing per-horizon metrics and feature importance.
\end{itemize}

The frontend is deployed on \textbf{Vercel}, with API calls proxied to the FastAPI Cloud backend via \texttt{vercel.json} rewrites.


% ══════════════════════════════════════════════════════════════════
\section{Automation (CI/CD)}

Pipelines are automated via \textbf{GitHub Actions} (\texttt{.github/workflows/pipelines.yml}):

\begin{itemize}
  \item \textbf{Feature pipeline}: Runs \textbf{every hour} (cron \texttt{0 * * * *}).  Fetches the last 14 days of actuals plus 3 days of forecasts, engineers features, and upserts them into Hopsworks.  Uses a minimal dependency set (\texttt{requirements-feature.txt} --- no PyTorch/XGBoost) to keep CI fast.
  \item \textbf{Training pipeline}: Runs \textbf{once daily at 02:00 UTC} (cron \texttt{0 2 * * *}).  Retrains all XGBoost and Random Forest models on the latest data (LSTM is disabled in CI to keep runtime under 10 minutes).  After training, updated model artifacts are committed back to the repository, which triggers a FastAPI Cloud auto-redeploy.
  \item \textbf{Manual dispatch}: Both pipelines can be triggered manually from the GitHub UI, with an option to run training alongside the feature update.
\end{itemize}

The training pipeline depends on the feature pipeline completing first, ensuring the feature store has the latest data before retraining begins.


% ══════════════════════════════════════════════════════════════════
\section{Key Design Decisions}

\begin{enumerate}[label=\textbf{\arabic*.}]
  \item \textbf{Open-Meteo over AQICN}: Open-Meteo provides both weather and AQI from a single API, includes forecast data, and requires no API key --- simplifying both development and CI/CD secrets management.

  \item \textbf{Direct vs.\ recursive forecasting}: The direct approach (separate model per horizon) avoids error accumulation at long horizons and allows horizon-specific feature selection.

  \item \textbf{Dropping raw pollutants from features}: Since US AQI is a deterministic function of pollutant concentrations, including current-timestep pollutant values as features would constitute target leakage.  Only \emph{lagged} and \emph{rolling} derivatives are retained.

  \item \textbf{Cyclical time encoding}: Sin/cos encoding for hour, month, and day-of-week ensures temporal proximity is preserved (e.g., hour 23 is treated as close to hour 0, not far away).

  \item \textbf{Separate CI requirements files}: Three \texttt{requirements-*.txt} files (feature, training, full) keep CI runners fast by only installing what each pipeline needs.  The feature pipeline avoids installing PyTorch entirely.

  \item \textbf{LSTM disabled in daily CI}: LSTM training takes $\sim$45 minutes per horizon and offers no advantage over XGBoost in this dataset.  It runs only locally when needed.

  \item \textbf{XGBoost as the production model}: XGBoost consistently outperformed both Random Forest and LSTM across all horizons while being fast to train and small to deploy.

  \item \textbf{Auto-commit model artifacts}: GitHub Actions commits updated models to the repo after training, triggering FastAPI Cloud to redeploy --- creating a fully automated retrain-deploy loop with no manual intervention.
\end{enumerate}


% ══════════════════════════════════════════════════════════════════
\section{Limitations and Future Work}

\subsection{Current Limitations}

\begin{itemize}
  \item \textbf{Long-horizon accuracy}: Performance at +48h and +72h is limited ($R^2 \approx 0.13$--$0.20$).  This is an inherent difficulty of multi-day air quality forecasting, not a modeling failure.
  \item \textbf{Single city}: The system currently supports only Karachi.  Extending to multiple cities would require per-city models or a location-aware architecture.
  \item \textbf{No TensorFlow}: TensorFlow was listed in the project requirements but was excluded due to a protobuf version conflict with the Hopsworks SDK on Python 3.13.  PyTorch was used for the deep learning component instead.
  \item \textbf{No hyperparameter tuning automation}: The current hyperparameters were manually tuned and are hardcoded.  An automated tuning step (e.g., Optuna) could improve performance.
\end{itemize}

\subsection{Possible Extensions}

\begin{itemize}
  \item Integrate satellite-based aerosol observations (e.g., from Copernicus/CAMS) for additional input features.
  \item Add push notification alerts when forecasted AQI exceeds hazardous thresholds.
  \item Implement conformal prediction intervals to communicate forecast uncertainty to users.
  \item Extend to a multi-city deployment with a shared model architecture.
\end{itemize}


% ══════════════════════════════════════════════════════════════════
\section{Technology Stack Summary}

\begin{table}[H]
\centering
\caption{Technologies used in the project.}
\begin{tabular}{ll}
\toprule
\textbf{Component} & \textbf{Technology} \\
\midrule
Language            & Python 3.12 \\
ML Models           & XGBoost, scikit-learn (Random Forest), PyTorch (LSTM) \\
Feature Store       & Hopsworks (cloud, free tier) \\
Model Registry      & Hopsworks Model Registry \\
Data API            & Open-Meteo (weather + air quality) \\
Backend API         & FastAPI \\
Frontend            & React, TypeScript, Vite, Tailwind CSS v4, Recharts \\
CI/CD               & GitHub Actions \\
Backend Hosting     & FastAPI Cloud \\
Frontend Hosting    & Vercel \\
Explainability      & SHAP \\
EDA                 & Matplotlib, Seaborn, statsmodels \\
Version Control     & Git / GitHub \\
\bottomrule
\end{tabular}
\end{table}


% ══════════════════════════════════════════════════════════════════
\section{Conclusion}

This project demonstrates a complete, end-to-end machine learning pipeline for air quality forecasting --- from raw API data to a live dashboard that anyone can use.  The key technical contributions are the leakage-free feature engineering pipeline, the direct multi-horizon modeling strategy with forecast weather integration, and the fully automated train-deploy loop via GitHub Actions.

The system achieves strong accuracy at short horizons (RMSE~$<$~1 for +1h) and useful-but-limited accuracy at longer horizons ($R^2 = 0.51$ at +24h), reflecting the fundamental difficulty of multi-day air quality prediction.  SHAP analysis confirms the model relies on physically sensible features --- recent AQI history for short-term and seasonal/weather patterns for long-term --- and is free of target leakage.

\end{document}
